\documentclass[12pt,a4paper]{article}
\usepackage[utf8]{inputenc}
\usepackage[T1]{fontenc}
\usepackage[spanish]{babel}
\usepackage{geometry}
\usepackage{xcolor}
\usepackage{enumitem}
\usepackage{booktabs}
\usepackage{array}
\usepackage{fancyhdr}
\usepackage{titlesec}
\usepackage{tcolorbox}
\usepackage{mdframed}
\usepackage{amssymb}

\geometry{margin=2.5cm}

\definecolor{accessred}{RGB}{175,0,0}
\definecolor{darkblue}{RGB}{0,70,127}
\definecolor{lightgray}{RGB}{245,245,245}
\definecolor{noteblue}{RGB}{230,242,255}
\definecolor{exgreen}{RGB}{0,120,0}
\definecolor{warnyellow}{RGB}{255,248,220}
\definecolor{practicegreen}{RGB}{235,250,235}

\pagestyle{fancy}
\fancyhf{}
\fancyhead[L]{\textcolor{accessred}{\textbf{UT6 -- Ejercicios: Microsoft Access}}}
\fancyhead[R]{\textcolor{gray}{\thepage}}
\fancyfoot[C]{\textcolor{gray}{\small AOF -- Natalia Marcos Iglesias}}
\renewcommand{\headrulewidth}{1.5pt}
\renewcommand{\headrule}{\hbox to\headwidth{\color{accessred}\leaders\hrule height \headrulewidth\hfill}}

\titleformat{\section}{\Large\bfseries\color{accessred}}{}{0em}{Sesión \thesection\ -- }[\titlerule]
\titleformat{\subsection}{\large\bfseries\color{darkblue}}{}{0em}{\thesubsection\ }

\tcbuselibrary{skins,breakable}

\newtcolorbox{ejercicio}[2][]{
  colback=lightgray, colframe=darkblue,
  title={\textbf{Ejercicio #2}},
  fonttitle=\bfseries\color{white},
  colbacktitle=darkblue,
  breakable, enhanced, #1
}

\newtcolorbox{practica}[1]{
  colback=practicegreen, colframe=exgreen,
  title={\textbf{Práctica: #1}},
  fonttitle=\bfseries\color{white},
  colbacktitle=exgreen,
  breakable, enhanced
}

\newtcolorbox{reto}{
  colback=warnyellow, colframe=accessred,
  title={\textbf{$\star$ Reto extra}},
  fonttitle=\bfseries\color{white},
  colbacktitle=accessred,
  breakable, enhanced
}

\newtcolorbox{reflexion}{
  colback=noteblue, colframe=darkblue!60,
  title={\textbf{Reflexión / Debate}},
  fonttitle=\bfseries,
  breakable, enhanced
}

\newcommand{\tecla}[1]{\fbox{\small\texttt{#1}}}

\title{
  \vspace{-1cm}
  {\color{accessred}\rule{\linewidth}{3pt}}\\[0.3cm]
  {\Huge\bfseries\color{accessred} UT6: Ejercicios Prácticos\\[0.2cm]
  Microsoft Access}\\[0.3cm]
  {\color{accessred}\rule{\linewidth}{1pt}}\\[0.2cm]
  {\large\color{gray} Módulo: Aplicaciones Ofimáticas (AOF)}
}
\author{\textbf{Natalia Marcos Iglesias}}
\date{}

\begin{document}

\maketitle
\thispagestyle{empty}

\begin{tcolorbox}[colback=noteblue, colframe=darkblue, title=\textbf{Base de datos del curso: ``Tienda Informática''}]
A lo largo de todas las sesiones trabajaremos con la base de datos \textbf{TiendaInformatica.accdb}.\\
Esta base de datos modela una tienda de componentes informáticos y contiene las siguientes tablas:

\begin{itemize}
  \item \textbf{Clientes}: CodCliente (PK), Nombre, Apellidos, Ciudad, Teléfono, Email, FechaNacimiento
  \item \textbf{Productos}: CodProducto (PK), Nombre, Categoria, PrecioUnidad, Stock
  \item \textbf{Proveedores}: CodProveedor (PK), NombreEmpresa, Ciudad, Teléfono
  \item \textbf{Pedidos}: CodPedido (PK), CodCliente (FK), FechaPedido, Enviado (Sí/No)
  \item \textbf{DetallesPedido}: CodDetalle (PK), CodPedido (FK), CodProducto (FK), Cantidad, Descuento
\end{itemize}
\end{tcolorbox}

\tableofcontents
\newpage

% ============================
\section{Conceptos básicos e interfaz de Access}
% ============================

\subsection{Ejercicios teóricos}

\begin{ejercicio}{1.1 -- Vocabulario de bases de datos}
Relaciona cada término con su definición correcta:

\begin{center}
\renewcommand{\arraystretch}{1.4}
\begin{tabular}{p{4cm} | p{7cm}}
\textbf{Término} & \textbf{Definición} \\
\hline
Campo & a) Conjunto de todos los campos de un mismo elemento \\
Registro & b) Identifica de forma única cada fila de una tabla \\
Clave primaria & c) Columna que almacena un tipo de dato \\
Clave foránea & d) Sistema de gestión de bases de datos \\
SGBD & e) Columna que referencia a la clave primaria de otra tabla \\
\end{tabular}
\end{center}
\end{ejercicio}

\begin{ejercicio}{1.2 -- Identificar elementos}
En la siguiente tabla de \textbf{CLIENTES}, señala:
\begin{enumerate}[label=\alph*)]
  \item ¿Cuántos registros tiene?
  \item ¿Cuántos campos tiene?
  \item ¿Qué campo sería la clave primaria más adecuada? ¿Por qué?
  \item ¿Podría el campo ``Nombre'' ser clave primaria? Justifica.
\end{enumerate}

\begin{center}
\renewcommand{\arraystretch}{1.2}
\begin{tabular}{|c|c|c|c|c|}
\hline
\textbf{CodCliente} & \textbf{Nombre} & \textbf{Apellidos} & \textbf{Ciudad} & \textbf{Teléfono} \\
\hline
C001 & Ana & García López & Valencia & 612345678 \\
C002 & Pedro & Martínez Gil & Madrid & 698765432 \\
C003 & María & García López & Sevilla & 611223344 \\
C004 & Luisa & Fernández Díaz & Valencia & 677889900 \\
\hline
\end{tabular}
\end{center}
\end{ejercicio}

\begin{ejercicio}{1.3 -- Tipos de relaciones}
Clasifica cada caso como relación 1:1, 1:N o N:M y explica por qué:
\begin{enumerate}[label=\alph*)]
  \item Personas -- DNI
  \item Profesores -- Asignaturas
  \item Estudiantes -- Cursos (un estudiante puede matricularse en varios cursos y un curso puede tener varios estudiantes)
  \item Países -- Capital
  \item Médicos -- Pacientes
\end{enumerate}
\end{ejercicio}

\begin{reflexion}
¿Para qué crees que sirve una base de datos frente a una simple hoja de cálculo de Excel? ¿Qué ventajas e inconvenientes tiene cada opción? Debatid en grupos de 3 personas durante 5 minutos.
\end{reflexion}

\subsection{Ejercicios prácticos}

\begin{practica}{Explorar la interfaz de Access}
\begin{enumerate}
  \item Abre Microsoft Access y describe en tu cuaderno los elementos que ves en la pantalla inicial.
  \item Crea una nueva base de datos en blanco y guárdala como \textbf{TiendaInformatica.accdb} en tu carpeta de trabajo.
  \item Identifica y señala en pantalla: Barra de acceso rápido, Cinta de opciones, Panel de navegación, Barra de estado.
  \item Practica contraer y expandir la cinta de opciones y el panel de navegación.
  \item Cierra la base de datos sin cerrar Access. Vuelve a abrirla desde \menu{Recientes}.
\end{enumerate}
\end{practica}

% ============================
\section{Creación de tablas}
% ============================

\subsection{Ejercicios teóricos}

\begin{ejercicio}{2.1 -- Tipos de datos}
Para cada campo, indica el tipo de dato más adecuado de Access y justifica tu elección:

\begin{center}
\renewcommand{\arraystretch}{1.4}
\begin{tabular}{p{5cm} | p{4cm} | p{4cm}}
\textbf{Campo} & \textbf{Tipo de dato} & \textbf{Justificación} \\
\hline
DNI del cliente & & \\
Precio de un producto & & \\
Fecha de nacimiento & & \\
¿Está pagada la factura? & & \\
Descripción larga de producto & & \\
Número de empleado (automático) & & \\
Foto del empleado & & \\
Código postal & & \\
Número de unidades en stock & & \\
\end{tabular}
\end{center}
\end{ejercicio}

\begin{ejercicio}{2.2 -- Diseñar la estructura de una tabla}
Diseña la estructura de la tabla \textbf{Empleados} para una empresa. Indica para cada campo: nombre, tipo de dato, tamaño y si sería clave primaria. La tabla debe almacenar: código de empleado, nombre, apellidos, fecha de nacimiento, cargo, salario mensual, foto, activo (trabaja actualmente o no), departamento.
\end{ejercicio}

\subsection{Ejercicios prácticos}

\begin{practica}{Crear la tabla Clientes}
Crea en Vista Diseño la tabla \textbf{Clientes} con los siguientes campos:

\begin{center}
\renewcommand{\arraystretch}{1.3}
\begin{tabular}{|l|l|l|}
\hline
\textbf{Campo} & \textbf{Tipo} & \textbf{Propiedades} \\
\hline
CodCliente & Texto corto & Tamaño: 5, Clave primaria \\
Nombre & Texto corto & Tamaño: 30, Requerido: Sí \\
Apellidos & Texto corto & Tamaño: 60, Requerido: Sí \\
Ciudad & Texto corto & Tamaño: 40, Valor predeterminado: ``Valencia'' \\
Teléfono & Texto corto & Tamaño: 15, Máscara: teléfono \\
Email & Hipervínculo & \\
FechaNacimiento & Fecha/Hora & Formato: Fecha corta \\
\hline
\end{tabular}
\end{center}

\begin{enumerate}[label=\alph*)]
  \item Asigna la clave principal al campo CodCliente.
  \item Guarda la tabla como \textbf{Clientes}.
  \item Cambia a Vista Hoja de datos e introduce 5 registros de clientes inventados.
\end{enumerate}
\end{practica}

\begin{practica}{Crear las tablas Productos y Pedidos}
\begin{enumerate}
  \item Crea la tabla \textbf{Productos} con: CodProducto (PK, Autonumeración), Nombre (Texto, 100), Categoria (Texto, 30), PrecioUnidad (Moneda), Stock (Número, Entero).
  \item Crea la tabla \textbf{Pedidos} con: CodPedido (PK, Autonumeración), CodCliente (Texto, 5), FechaPedido (Fecha/Hora), Enviado (Sí/No).
  \item Añade una regla de validación al campo Stock: el valor no puede ser negativo. Escribe también el texto de validación apropiado.
  \item Introduce al menos 8 productos y 5 pedidos.
\end{enumerate}
\end{practica}

\begin{reto}
Añade a la tabla \textbf{Clientes} un campo calculado llamado \textbf{Edad} que muestre la edad actual del cliente usando la función \texttt{Year(Date()) - Year([FechaNacimiento])}. Comprueba que funciona correctamente en Vista Hoja de datos.
\end{reto}

% ============================
\section{Modificar tablas y propiedades}
% ============================

\begin{ejercicio}{3.1 -- Propiedades de campo}
Indica qué propiedad de campo utilizarías en cada caso:
\begin{enumerate}[label=\alph*)]
  \item Quieres que el campo ``País'' muestre ``España'' automáticamente.
  \item Necesitas que el campo ``Salario'' solo acepte valores entre 1.000 y 10.000.
  \item Quieres que el campo ``NIF'' muestre la plantilla de entrada \texttt{00000000-X}.
  \item Necesitas que el campo ``FechaAlta'' sea obligatorio.
  \item Quieres que el campo ``NúmEmpleado'' aparezca como encabezado ``Número de empleado''.
\end{enumerate}
\end{ejercicio}

\begin{practica}{Modificar la tabla Clientes}
\begin{enumerate}
  \item Abre la tabla \textbf{Clientes} en Vista Diseño.
  \item Añade el campo \textbf{Descuento} (Número, Simple) con valor predeterminado 0 y regla de validación \texttt{>=0 Y <=50}, con el texto ``El descuento debe estar entre 0 y 50\%''.
  \item Añade el campo \textbf{Activo} (Sí/No) con valor predeterminado Sí.
  \item Cambia el título del campo \textbf{FechaNacimiento} a ``Fecha de Nacimiento''.
  \item Cambia el índice del campo \textbf{Apellidos} a ``Sí (con duplicados)''.
  \item Guarda los cambios e introduce datos en los nuevos campos.
\end{enumerate}
\end{practica}

\begin{practica}{Buscar y reemplazar datos}
\begin{enumerate}
  \item En la tabla \textbf{Clientes}, usa \tecla{Ctrl+B} para buscar todos los clientes de ``Valencia''.
  \item Usa la función \textbf{Reemplazar} para cambiar la ciudad ``Madrid'' por ``Madrid (capital)''.
  \item Desplázate hasta el último registro, luego al primero, usando los botones de navegación.
\end{enumerate}
\end{practica}

% ============================
\section{Relaciones entre tablas}
% ============================

\begin{ejercicio}{4.1 -- Diseño de relaciones}
Para la base de datos \textbf{TiendaInformatica}, dibuja el esquema de relaciones indicando:
\begin{itemize}
  \item Nombre de cada tabla y sus campos.
  \item Las relaciones entre tablas (qué campos se relacionan).
  \item El tipo de cada relación (1:1, 1:N, N:M).
\end{itemize}
\end{ejercicio}

\begin{practica}{Crear las relaciones}
\begin{enumerate}
  \item Crea también la tabla \textbf{DetallesPedido}: CodDetalle (PK, Autonumeración), CodPedido (Número, Entero largo), CodProducto (Número, Entero largo), Cantidad (Número, Entero), Descuento (Número, Simple).
  \item Abre la ventana de Relaciones (\menu{Herramientas de base de datos > Relaciones}).
  \item Crea las siguientes relaciones con integridad referencial:
    \begin{itemize}
      \item \textbf{Clientes.CodCliente} $\leftrightarrow$ \textbf{Pedidos.CodCliente} (1:N)
      \item \textbf{Pedidos.CodPedido} $\leftrightarrow$ \textbf{DetallesPedido.CodPedido} (1:N)
      \item \textbf{Productos.CodProducto} $\leftrightarrow$ \textbf{DetallesPedido.CodProducto} (1:N)
    \end{itemize}
  \item Activa la opción ``Actualizar en cascada'' en todas las relaciones.
  \item Guarda las relaciones e introduce datos coherentes en DetallesPedido.
\end{enumerate}
\end{practica}

\begin{ejercicio}{4.2 -- Integridad referencial}
Responde razonadamente:
\begin{enumerate}[label=\alph*)]
  \item ¿Qué ocurriría si intentas insertar un pedido con un CodCliente que no existe en la tabla Clientes? ¿Y si no hay integridad referencial?
  \item ¿Qué diferencia hay entre ``Actualizar en cascada'' y ``Eliminar en cascada''?
  \item ¿Por qué no es buena idea poner ``Eliminar en cascada'' en todas las relaciones? Pon un ejemplo de cuándo podría ser peligroso.
\end{enumerate}
\end{ejercicio}

% ============================
\section{Consultas básicas}
% ============================

\begin{practica}{Consultas de selección simples}
Crea las siguientes consultas de selección y guárdalas con el nombre indicado:
\begin{enumerate}
  \item \textbf{ConsClientesValencia}: Muestra nombre, apellidos y teléfono de todos los clientes de Valencia, ordenados alfabéticamente por apellidos.
  \item \textbf{ConsProductosCara}: Muestra todos los datos de productos con precio superior a 100€, ordenados de mayor a menor precio.
  \item \textbf{ConsPedidosPendientes}: Muestra los pedidos no enviados (Enviado = No) con su fecha, ordenados por fecha de más antiguo a más reciente.
  \item \textbf{ConsClientesMayores}: Muestra clientes nacidos antes del 1/1/1990 con descuento mayor del 10\%.
\end{enumerate}
\end{practica}

\begin{practica}{Consultas con campos calculados}
\begin{enumerate}
  \item Crea la consulta \textbf{ConsPrecioIVA}: muestra CodProducto, Nombre, PrecioUnidad y un campo calculado \textbf{Precio con IVA} (PrecioUnidad * 1,21), con el encabezado ``Precio c/IVA''.
  \item Crea la consulta \textbf{ConsValorStock}: muestra Nombre del producto, Stock y un campo calculado \textbf{ValorStock} (Stock * PrecioUnidad) que indica el valor total del stock de cada producto.
  \item Formatea el campo ValorStock como Moneda.
\end{enumerate}
\end{practica}

\begin{practica}{Consultas con comodines}
Crea las siguientes consultas usando el operador \textbf{Como}:
\begin{enumerate}
  \item \textbf{ConsNombresA}: Clientes cuyo nombre empiece por la letra ``A''.
  \item \textbf{ConsEmailGoogle}: Clientes con email de Gmail (que contenga ``@gmail'').
  \item \textbf{ConsProductosRAM}: Productos cuyo nombre contenga la palabra ``RAM''.
\end{enumerate}
\end{practica}

\begin{reto}
Crea la consulta \textbf{ConsClientesCumpleanos}: muestra los clientes que cumplen años este mes. Usa la función \texttt{Month([FechaNacimiento]) = Month(Date())}. ¡Puede ser útil para una campaña de marketing!
\end{reto}

% ============================
\section{Consultas avanzadas}
% ============================

\begin{practica}{Consultas con parámetros}
\begin{enumerate}
  \item Crea la consulta \textbf{ConsBuscarCiudad}: pide al usuario una ciudad y muestra los clientes de esa ciudad.
  \item Crea la consulta \textbf{ConsBuscarPrecio}: pide al usuario un precio máximo y muestra los productos más baratos que ese precio.
  \item Crea la consulta \textbf{ConsBuscarFechas}: pide fecha de inicio y fecha de fin, y muestra los pedidos en ese rango de fechas.
\end{enumerate}
\end{practica}

\begin{practica}{Consultas multitabla}
\begin{enumerate}
  \item Crea la consulta \textbf{ConsClientesPedidos}: muestra nombre y apellidos del cliente, fecha del pedido y si está enviado.
  \item Crea la consulta \textbf{ConsDetallesCompletos}: muestra para cada línea de pedido el nombre del cliente, la fecha del pedido, el nombre del producto, la cantidad y el descuento.
  \item Modifica la consulta anterior para mostrar también el subtotal de cada línea: \texttt{Cantidad * PrecioUnidad * (1 - Descuento/100)}.
\end{enumerate}
\end{practica}

\begin{ejercicio}{6.1 -- Composición externa}
\begin{enumerate}[label=\alph*)]
  \item ¿Cuál es la diferencia entre una composición interna y una externa?
  \item Crea una consulta que muestre TODOS los clientes, con sus pedidos si los tienen, y NULL si no han hecho ningún pedido. (Pista: composición externa izquierda).
  \item ¿Para qué sería útil esta consulta en una empresa real?
\end{enumerate}
\end{ejercicio}

\begin{reto}
Crea la consulta \textbf{ConsClientesSinPedidos}: muestra los clientes que NO han realizado ningún pedido. Usa una composición externa y el criterio ``Es Nulo'' en el campo CodPedido.
\end{reto}

% ============================
\section{Consultas de resumen y referencias cruzadas}
% ============================

\begin{practica}{Consultas de resumen}
Crea las siguientes consultas de resumen:
\begin{enumerate}
  \item \textbf{ConsResumenCiudad}: Número de clientes por ciudad, ordenado descendentemente.
  \item \textbf{ConsEstadisticaProductos}: Para cada categoría de producto, muestra el precio mínimo, máximo, promedio y el número de productos.
  \item \textbf{ConsVentasTotales}: Suma del importe total de cada pedido (usando DetallesPedido: Cantidad * PrecioUnidad).
  \item \textbf{ConsTop5Productos}: Los 5 productos más vendidos (por cantidad total). Activa la propiedad ``Valores superiores'' en el diseño de la consulta.
\end{enumerate}
\end{practica}

\begin{practica}{Consulta de referencias cruzadas}
Crea con el asistente la consulta \textbf{CruzVentasMensuales}:
\begin{itemize}
  \item Filas: Año del pedido.
  \item Columnas: Mes del pedido.
  \item Valor: Suma de importes (Cantidad * PrecioUnidad).
\end{itemize}
Esta consulta permitirá ver la evolución de ventas mes a mes.
\end{practica}

\begin{reflexion}
¿Qué información de negocio podrías extraer de estas consultas de resumen? ¿Qué decisiones comerciales podrías tomar? Piensa en 3 ejemplos concretos y debátelos con tu compañero.
\end{reflexion}

% ============================
\section{Consultas de acción}
% ============================

\begin{importante}
\textbf{Antes de ejecutar cualquier consulta de acción}, usa la Vista Hoja de datos para revisar qué registros se verán afectados. Las consultas de acción son \textbf{irreversibles}.
\end{importante}

\begin{practica}{Consultas de acción}
\begin{enumerate}
  \item \textbf{Actualización}: Crea la consulta \textbf{SubirPrecio10}: sube un 10\% el precio de todos los productos de la categoría ``RAM''. Antes de ejecutar, revisa qué productos afecta.

  \item \textbf{Creación de tabla}: Crea la consulta \textbf{CrearClientesVIP}: genera una nueva tabla llamada \textbf{ClientesVIP} con los clientes que tienen descuento mayor del 20\%.

  \item \textbf{Datos anexados}: Crea la consulta \textbf{AnexarClientesPrueba}: añade a \textbf{ClientesVIP} los clientes con más de 3 pedidos realizados. (Tendrás que crear primero una consulta de resumen que cuente los pedidos por cliente.)

  \item \textbf{Eliminación}: Crea la consulta \textbf{EliminarPedidosAntiguos}: elimina los pedidos con fecha anterior al 01/01/2020 que ya estén marcados como enviados.
\end{enumerate}
\end{practica}

\begin{ejercicio}{8.1 -- Consultas de acción}
Responde:
\begin{enumerate}[label=\alph*)]
  \item ¿Qué diferencia hay entre la consulta de creación de tabla y la de datos anexados?
  \item ¿Qué pasa con la clave principal cuando creas una tabla nueva con una consulta de creación?
  \item ¿Por qué es peligroso una consulta de eliminación sin criterio de búsqueda?
  \item Si la tabla de destino tiene activada la integridad referencial, ¿qué restricciones aplican a la consulta de eliminación?
\end{enumerate}
\end{ejercicio}

% ============================
\section{Formularios}
% ============================

\begin{practica}{Crear formularios con el asistente}
\begin{enumerate}
  \item Crea con el asistente el formulario \textbf{FrmClientes} basado en la tabla Clientes. Usa distribución \emph{En columnas} e incluye todos los campos.
  \item Crea el formulario \textbf{FrmProductos} basado en la tabla Productos, distribución \emph{Tabular}.
  \item Usa los formularios para introducir 5 clientes y 5 productos nuevos.
  \item Practica navegar entre registros con los botones de la barra inferior.
\end{enumerate}
\end{practica}

\begin{practica}{Personalizar un formulario en Vista Diseño}
Modifica el formulario \textbf{FrmClientes}:
\begin{enumerate}
  \item Añade en el encabezado una etiqueta con el texto ``GESTIÓN DE CLIENTES'' en negrita, tamaño 16.
  \item Cambia el color de fondo del formulario.
  \item Reorganiza los controles: pon Nombre y Apellidos en la misma fila.
  \item Añade un botón de comando para cerrar el formulario (usa el asistente del botón).
  \item Cambia la propiedad \textbf{Título} del formulario a ``Ficha de Cliente''.
\end{enumerate}
\end{practica}

\begin{practica}{Formulario con subformulario}
\begin{enumerate}
  \item Crea el formulario \textbf{FrmPedidosCompleto}:
    \begin{itemize}
      \item Formulario principal basado en \textbf{Pedidos}: muestra los datos del pedido y del cliente.
      \item Subformulario basado en \textbf{DetallesPedido}: muestra las líneas del pedido.
    \end{itemize}
  \item Verifica que al navegar entre pedidos el subformulario muestra correctamente los detalles de cada uno.
  \item Añade en el pie del subformulario un campo calculado que muestre el \textbf{total del pedido}.
\end{enumerate}
\end{practica}

\begin{reto}
Crea un formulario de \textbf{búsqueda de clientes}: con un cuadro de texto donde el usuario escriba un nombre o ciudad, y un botón que ejecute una consulta con parámetro y muestre los resultados en un subformulario. (Pista: usa macros o VBA para ejecutar la consulta al pulsar el botón.)
\end{reto}

% ============================
\section{Proyecto integrador final}
% ============================

\begin{tcolorbox}[colback=lightgray, colframe=accessred, title=\textbf{Proyecto Final: Base de datos completa}, colbacktitle=accessred, fonttitle=\bfseries\color{white}]
Usando todo lo aprendido, amplía y mejora la base de datos \textbf{TiendaInformatica.accdb} con los siguientes elementos adicionales:
\end{tcolorbox}

\begin{practica}{Proyecto final -- parte 1: Tablas y relaciones}
\begin{enumerate}
  \item Añade la tabla \textbf{Empleados}: CodEmpleado (PK), Nombre, Apellidos, Cargo, Salario, FechaAlta.
  \item Añade la tabla \textbf{Proveedores}: CodProveedor (PK), NombreEmpresa, Ciudad, Teléfono.
  \item Añade el campo \textbf{CodProveedor} a la tabla Productos para saber qué proveedor suministra cada producto.
  \item Añade el campo \textbf{CodEmpleado} a la tabla Pedidos para saber qué empleado gestionó cada pedido.
  \item Crea todas las relaciones necesarias con integridad referencial.
\end{enumerate}
\end{practica}

\begin{practica}{Proyecto final -- parte 2: Consultas}
Crea las siguientes consultas:
\begin{enumerate}
  \item \textbf{ConsInformeVentas}: muestra nombre del empleado, número de pedidos gestionados y suma total de ventas.
  \item \textbf{ConsProductosBajoStock}: productos con stock inferior a 10 unidades, ordenados por categoría.
  \item \textbf{ConsClientesMasActivos}: los 10 clientes que más pedidos han realizado.
  \item \textbf{ConsFacturacion}: muestra para cada pedido el nombre del cliente, empleado responsable, fecha, número de líneas y total del pedido.
\end{enumerate}
\end{practica}

\begin{practica}{Proyecto final -- parte 3: Formularios}
\begin{enumerate}
  \item Crea un formulario de \textbf{menú principal} con botones para abrir cada formulario de la aplicación.
  \item Crea formularios completos para gestionar Empleados y Proveedores.
  \item Personaliza el aspecto de todos los formularios con colores, fuentes y un logo de la empresa (puedes usar cualquier imagen).
\end{enumerate}
\end{practica}

\begin{reflexion}
\textbf{Presentación final (5 min por grupo):} Muestra al resto de la clase tu base de datos terminada. Explica las decisiones que tomaste en el diseño (tipos de datos, relaciones, consultas útiles) y qué mejorarías si tuvieras más tiempo.
\end{reflexion}

\vfill
\begin{center}
{\color{accessred}\rule{0.5\linewidth}{1pt}}\\[0.3cm]
{\small\color{gray} UT6 -- Ejercicios prácticos de Microsoft Access \quad|\quad AOF}
\end{center}

\end{document}
